% GENERATED FILE. Edit canonical lessons, then run npm run generate:books.
% canonical-source-hash: fnv1a64:356e4d1c5d90dbf1
% canonical-lessons: RU-C187-polnyy, RU-C187-glubokiy, RU-C187-rovnyy, RU-C187-kvadratnyy, RU-C187-ostryy

\chapter{Full, Deep, Even, Square, Sharp}
\label{ch:ru-full-deep-even-square-sharp}
\hlchaptermodality{187}

\begin{chapteropening}
\textbf{By the end of this chapter:} \emph{I can say full, deep, even, square and sharp.}

Complete the last lesson of chapter 187: I can say full, deep, even, square and sharp.
\end{chapteropening}

\section[\texorpdfstring{\ru{полный} (pólnyy)}{полный (pólnyy)}]{\ru{полный} (pólnyy) — full}

\label{lesson:RU-C187-polnyy}



\begin{quote}
Before the new one: say the Russian for generous, then the Russian for greedy.
\end{quote}

\subsection*{You'll want to know: \ru{полный}}
\textbf{\ru{полный}} (\emph{pólnyy}) — \textquotedblleft{}full\textquotedblright{}.

\begin{grammarlens}[title={What you've built}]
One more word for this chapter.
\end{grammarlens}

\subsection*{Your turn}
\begin{itemize}
\raggedright
  \item \emph{Say it:} \emph{pólnyy}
  \item \emph{Say it:} \emph{pólnyy}, once more
  \item \emph{Say it:} say \emph{pólnyy}
  \item \emph{Recall:} say the Russian for generous, then the Russian for greedy, then say \emph{pólnyy} again
\end{itemize}

\begin{tcolorbox}[breakable,colback=teal!4,colframe=teal!35!black,title={Before you move on}]
What does \ru{полный} mean? (\textquotedblleft{}Full\textquotedblright{}.) Say it once more.
\end{tcolorbox}
\section[\texorpdfstring{\ru{глубокий} (glubókiy)}{глубокий (glubókiy)}]{\ru{глубокий} (glubókiy) — deep}

\label{lesson:RU-C187-glubokiy}



\begin{quote}
Before the new one: say the Russian for greedy, then the Russian for full.
\end{quote}

\subsection*{You'll want to know: \ru{глубокий}}
\textbf{\ru{глубокий}} (\emph{glubókiy}) — \textquotedblleft{}deep\textquotedblright{}.

\begin{grammarlens}[title={What you've built}]
One more word for this chapter.
\end{grammarlens}

\subsection*{Your turn}
\begin{itemize}
\raggedright
  \item \emph{Say it:} \emph{glubókiy}
  \item \emph{Say it:} \emph{glubókiy}, once more
  \item \emph{Say it:} say \emph{glubókiy}
  \item \emph{Recall:} say the Russian for greedy, then the Russian for full, then say \emph{glubókiy} again
\end{itemize}

\begin{tcolorbox}[breakable,colback=teal!4,colframe=teal!35!black,title={Before you move on}]
What does \ru{глубокий} mean? (\textquotedblleft{}Deep\textquotedblright{}.) Say it once more.
\end{tcolorbox}
\section[\texorpdfstring{\ru{ровный} (róvnyy)}{ровный (róvnyy)}]{\ru{ровный} (róvnyy) — even, level}

\label{lesson:RU-C187-rovnyy}



\begin{quote}
Before the new one: say the Russian for full, then the Russian for deep.
\end{quote}

\subsection*{You'll want to know: \ru{ровный}}
\textbf{\ru{ровный}} (\emph{róvnyy}) — \textquotedblleft{}even, level\textquotedblright{}.

\begin{grammarlens}[title={What you've built}]
One more word for this chapter.
\end{grammarlens}

\subsection*{Your turn}
\begin{itemize}
\raggedright
  \item \emph{Say it:} \emph{róvnyy}
  \item \emph{Say it:} \emph{róvnyy}, once more
  \item \emph{Say it:} say \emph{róvnyy}
  \item \emph{Recall:} say the Russian for full, then the Russian for deep, then say \emph{róvnyy} again
\end{itemize}

\begin{tcolorbox}[breakable,colback=teal!4,colframe=teal!35!black,title={Before you move on}]
What does \ru{ровный} mean? (\textquotedblleft{}Even, level\textquotedblright{}.) Say it once more.
\end{tcolorbox}
\section[\texorpdfstring{\ru{квадратный} (kvadrátnyy)}{квадратный (kvadrátnyy)}]{\ru{квадратный} (kvadrátnyy) — square}

\label{lesson:RU-C187-kvadratnyy}



\begin{quote}
Before the new one: say the Russian for deep, then the Russian for even.
\end{quote}

\subsection*{You'll want to know: \ru{квадратный}}
\textbf{\ru{квадратный}} (\emph{kvadrátnyy}) — \textquotedblleft{}square\textquotedblright{}.

\begin{grammarlens}[title={What you've built}]
One more word for this chapter.
\end{grammarlens}

\subsection*{Your turn}
\begin{itemize}
\raggedright
  \item \emph{Say it:} \emph{kvadrátnyy}
  \item \emph{Say it:} \emph{kvadrátnyy}, once more
  \item \emph{Say it:} say \emph{kvadrátnyy}
  \item \emph{Recall:} say the Russian for deep, then the Russian for even, then say \emph{kvadrátnyy} again
\end{itemize}

\begin{tcolorbox}[breakable,colback=teal!4,colframe=teal!35!black,title={Before you move on}]
What does \ru{квадратный} mean? (\textquotedblleft{}Square\textquotedblright{}.) Say it once more.
\end{tcolorbox}
\section[\texorpdfstring{\ru{острый} (óstryy)}{острый (óstryy)}]{\ru{острый} (óstryy) — sharp; spicy}

\label{lesson:RU-C187-ostryy}



\begin{quote}
Before the new one: say the Russian for even, then the Russian for square.
\end{quote}

\subsection*{You'll want to know: \ru{острый}}
\textbf{\ru{острый}} (\emph{óstryy}) — \textquotedblleft{}sharp; spicy\textquotedblright{}.

\begin{grammarlens}[title={What you've built}]
One more word for this chapter.
\end{grammarlens}

\subsection*{Your turn}
\begin{itemize}
\raggedright
  \item \emph{Say it:} \emph{óstryy}
  \item \emph{Say it:} \emph{óstryy}, once more
  \item \emph{Say it:} say \emph{óstryy}
  \item \emph{Recall:} say the Russian for even, then the Russian for square, then say \emph{óstryy} again
\end{itemize}

\begin{tcolorbox}[breakable,colback=teal!4,colframe=teal!35!black,title={Before you move on}]
What does \ru{острый} mean? (\textquotedblleft{}Sharp; spicy\textquotedblright{}.) Say it once more.
\end{tcolorbox}
